%! TEX root = ../main.tex
\documentclass[main]{subfiles}

\begin{document}
\chapter{はじめに}
\label{chap:introduction}
% ここを書き換える: 研究の背景・目的・論文の構成
\section{背景}
この章では、研究の背景と目的を述べる。
まず、研究対象の分野における現状と課題を、先行研究を引用しながら説明する。
引用は \texttt{references.bib} に文献情報を登録し、\verb|\cite{キー}| で参照する\cite{lamport1994latex}。
複数の文献を同時に引用することもできる\cite{knuth1986texbook,knuth1981breaking}。
Webページは \texttt{@online} として登録し、URLと閲覧日を記載する\cite{ctan_luatexja}。\par

\section{目的}
背景で述べた課題に対し、本研究で何を明らかにするのか、あるいは何を実現するのかを明確に述べる。
本研究の貢献は次のとおりである。
\begin{itemize}
  \item 貢献1: 〇〇の課題に対する新しい手法を提案する
  \item 貢献2: 提案手法を△△のデータで評価し、有効性を示す
\end{itemize}

\section{本論文の構成}
本論文は次のように構成される。
\ref{chap:method}章では提案手法について述べる。
\ref{chap:experiment}章では評価実験の内容を述べ、\ref{chap:result}章で結果を示して考察を行う。
\ref{chap:conclusion}章では本論文の成果をまとめ、今後の課題を述べる。
\end{document}
